\section{Case: refinement proofs}
\label{sec:refinement}

\paragraph{Background.} When a proof assistant is used to verify an operating-system
kernel such as seL4 \citep{klein2009} or a compiler such as CompCert \citep{leroy2009},
the standard method is \emph{refinement}. An implementation refines a specification when
every behaviour of the implementation, as far as it is visible at the level of the
specification, is a behaviour the specification allows. The usual proof technique links
the two by a simulation relation
$R \subseteq S_{\mathrm{impl}} \times S_{\mathrm{spec}}$, such that every concrete step
is matched by an abstract step, or by none (a stuttering step), and $R$ is preserved.
Matching steps in this way establishes the inclusion for safety properties. Liveness
needs a further argument, since an implementation that stutters forever matches every
step and does nothing \citep{abadi1991}.

\subsection{What one relation does not say}

What a refinement proof carries from the specification to the implementation is what the
inclusion of behaviours preserves, and no more. A property of single executions that
holds of every behaviour of the specification holds of every behaviour of the
implementation; that is the point of the method. Properties that compare executions do
not carry over in general. \citet{clarkson2010} show that refinement, which never
invalidates a property of single traces, can invalidate a hyperproperty such as
noninterference: a system that satisfies a confidentiality policy can have a refinement
that violates it. The security literature knows this as the refinement paradox.

The practice of the field confirms it one property at a time. The functional-correctness
proof of seL4 makes no claim about information flow or about timing. Information-flow
enforcement was established by a further proof \citep{murray2013}, and time protection
was added as a mechanism of its own \citep{ge2019}. The semantic preservation theorem of
CompCert does not say that compiled code keeps the constant-time behaviour of its source;
that took a modified compiler and a separate proof \citep{barthe2020}. Each time a
further property was carried across the same boundary, and each time it needed a
preservation argument of its own. In the vocabulary of Section~\ref{sec:casts} these are
move lemmas, registered one property at a time.

Where verified systems fail in use, the failures sit where this account expects them.
\citet{fonseca2017} examined three formally verified distributed systems and found
sixteen bugs. None was in the protocols that the proofs covered. They lay in what the
proofs had to assume: the interface to unverified code, the specifications, and the
tools that checked the proofs. What failed was the transfer of a predicate. ``Verified'' is
established of a model, under stated assumptions and for a stated property, and is then
said of the artefact. That is the reuse that predicate locality (A2) restricts: a
predicate carried to another layer without showing that its truth conditions hold there.

\subsection{Five structures behind one relation}

A refinement relation is usually drawn as one vertical arrow between two levels. The five
structures of Section~\ref{sec:vocabulary} each correspond to a question that the
relation, taken alone, leaves open (Table~\ref{tab:five}).

\begin{table}[t]
\centering\small
\begin{tabularx}{\textwidth}{@{}P{2.1cm}P{3.3cm}L@{}}
\toprule
Structure & In verification & What a single refinement relation leaves undeclared \\
\midrule
S1 Substrate &
  Hardware: registers, caches, memory cells. &
  Causal paths that run through the substrate and are absent from the model. Spectre
  leaks through cache state that the instruction-set model does not contain, and leaves
  the modelled state untouched \citep{kocher2019}. Rowhammer differs in kind: a fault
  flips bits, and so changes the modelled state by a transition the model does not have
  \citep{kim2014}. \\
S2 Event abstraction &
  Operational semantics: small-step, big-step, traces. &
  Big-step semantics relates initial states to final states. Intermediate states and
  durations are abstracted away, so an interleaving between two concrete steps of one
  abstract step cannot be stated at the abstract level. \\
S3 Representation &
  Bytes in memory, the parsed syntax tree, the abstract value. &
  A functional-correctness proof of a parser establishes correspondence: the tree
  represents the bytes. It is silent on provenance, and a verified parser parses a
  forged message correctly. Provenance is a different property with proofs of its own,
  such as authentication in protocol verification. \\
S4 Epistemic access &
  Observational equivalence, noninterference; monitors and logs. &
  What an observer can infer is a property of sets of executions, which refinement does
  not preserve in general. And an observation is produced by a process with a latency of
  its own, which is no part of the transition observed. \\
S5 Normative &
  Security policies, access-control rules, contracts. &
  Encoding a prohibition as an obligation to omit keeps every verdict, but the type of a
  norm then no longer says what an omission means
  (Section~\ref{sec:deontic}). \\
\bottomrule
\end{tabularx}
\caption{Five structures, and the question each raises that a refinement relation alone
does not answer.}
\label{tab:five}
\end{table}

\subsection{A refinement relation has no duration}
\label{sec:duration}

On the mechanistic account of levels, relations between levels are constitutive and not
causal. The higher-level state is realised by the lower-level states at the same time,
and what looks like causation between levels is causation within a level together with
constitution \citep{craver2007}. A simulation relation is of this kind. It says which
abstract state a concrete state amounts to or, for data, which abstract value some bytes
represent. Nothing passes from the one to the other and no time elapses. Such an edge has
no $\Dt$. The duration is undefined, not zero, because duration is a function on causal
edges alone. This is a statement about the relation, not about the two models it links.
Each of them can represent time and intermediate states in its own transitions, to
whatever extent its semantics provides.

Two points follow, one inside verification and one outside it.

The first concerns atomic steps. A step that is atomic in the specification may be
realised by several concrete steps with time between them. At the abstract level that
interval does not exist. Whether anything can happen in it is a question about the
concrete model. A refinement proof can establish that the concrete steps together
implement the atomic operation, and where the concrete model contains the interleavings
and the interference that can occur between them, an attack in the window would show as
a behaviour the specification forbids, and the proof would fail. The gap opens where the
concrete model leaves them out: a sequential model of code that runs concurrently, or an
environment assumed to be passive. Time-of-check to time-of-use races and reentrancy
have this form. The check and the use are one step above and two below, and the proof
covers the window exactly as far as the concrete model does.

The second lies in modelling consequences. In a risk model, the arrows from a system
event to a data event to a business event also run upward, from a technical description
to a business one. They are not refinement. In the consequence chain of the taxonomy for
which LEM was developed \citep{kreinztlctc}, malware alters a payment-generation process,
a false payment instruction is generated, and funds are transferred. These are three
events, not three descriptions of one. Time elapses between them, and payment screening
can intervene between the second and the third. Drawn as levels of one stack, they invite
the constitutive reading, and with it the window in which monitors and circuit-breaker
controls act disappears from the model. This is why the type of an edge is determined per
edge, by the two tests of Section~\ref{sec:vocabulary}, and not by which pair of layers
the edge connects.

\subsection{The normative crossing: three deontic types}
\label{sec:deontic}

The link from a step of a trace to the norm it fulfils or violates is also a crossing,
and here the vocabulary of the norm decides what can be read off its type.

Deontic logic can define a prohibition as an obligation to omit:
$\mathrm{Forbid}(\varphi) \equiv \mathrm{Oblige}(\neg\varphi)$ \citep{vonwright1951}. The
reduction is correct and costs nothing in expressive power. Every act and every omission
receives the same verdict before and after it.

What the reduction gives up is a diagnostic property, and the claim made for it here is
a restricted one. Consider a norm about one act, a norm that applies to the case, and an
agent who does not perform the act. Call this an omission. It is not the stuttering step
of a refinement proof: a step can leave the abstract state unchanged while the
implementation does a good deal, including the act in question.
\begin{itemize}
  \item Under a prohibition of the act, the omission is fulfilment.
  \item Under an obligation to perform the act that is due and not yet discharged, the
        omission is a violation.
  \item Under a right, in the sense of a permission, it is neither. The right is simply
        not exercised.
\end{itemize}
With three types, the type of the norm settles the standing of the omission. With two it
does not. $\mathrm{Oblige}(\varphi)$ and $\mathrm{Oblige}(\neg\varphi)$ are both
obligations, and the same omission violates the one and fulfils the other. Only the
content of the norm tells them apart. That is the whole benefit: with a third type, one
distinction can be read off the type without inspecting the content. LEM keeps three
types for this reason. The choice is a modelling convention. It is not an argument that
normative systems need three primitive types.

The mechanised model (Appendix~\ref{app:deontic}) is correspondingly small. It has no
deadline, no condition of applicability and no state of discharge: an obligation in it
is an obligation that is applicable, due and undischarged.

Why ``due'' matters has a standard formal ground, which depends on the content of a norm
and not on its label. A norm that an event never occur is a safety property: a violation
shows in a finite prefix of the trace. A norm that an act eventually be performed is a
liveness property: no finite prefix violates it, because the act may still come
\citep{alpern1985}. A deadline turns the second into something a monitor can find
violated. Prohibitions are typically of the first form and obligations of the second,
so a policy language that keeps the two types apart will tend to keep two kinds of check
apart, but the type alone does not guarantee it.

The right meant here is a permission, which \citet{hohfeld1913} calls a privilege. It is
not his claim-right, the correlative of another party's duty.

The normative edge is not classificatory: amending a counting rule changes what an event
is counted as, while amending a norm changes what is required. And it carries no $\Dt$. A
deadline is not a duration on this edge. It bounds a separate causal process, which runs
from the event that starts the clock to the act that discharges the duty.

\paragraph{What the field does, and what the rule adds.} Verification already proves the
preservation of each further property separately, and the theory of hyperproperties and
of secure compilation classifies what a refinement or a compiler preserves by class of
property (Section~\ref{sec:related}). The rule adds two things to the preservation list.
The first is the kind of the crossing: a simulation relation is constitutive or
representational and has no duration, so a sequence of events with time between them,
such as a consequence chain, should not be modelled as if it were a refinement. The
second is the observation that ``verified'' is itself a predicate with a layer: it is
established of a model, under that model's assumptions about concurrency and
environment, and needs an argument of its own to be said of an artefact.
